\subsection{Redistribution without aggregate inequality change}
The independent-transaction controls produce zero differences in measured allocation outcomes across all three selectors at every capacity. In four-slot chain cases, Core minus scalar unrelated-burden contrast is 0.583 blocks (95\% paired-bootstrap interval [0.508, 0.658]), or 350.0 scheduled seconds. First-block fees increase by 8,369.500 sat on average, yet full-cohort mean burden and Gini remain exactly unchanged because each selector still confirms four transactions in each of three blocks. The result is redistribution of which identities wait, as Proposition~1 predicts.

For the child-pays-for-parent pair, the corresponding unrelated contrast is 0.258 blocks [0.208, 0.308]. The mixed graph yields the same aggregate contrasts as the chain for these fee assignments; this does not establish their equivalence under general demand. At six slots, the chain contrast is 0.283 blocks [0.225, 0.342]. Table~\ref{tab:contrasts} and Fig.~\ref{fig:tradeoffs} retain every tested capacity and graph, including null differences. First-block utilization equals one for every policy and cell; capacity stress effects arise in subsequent blocks, not initial weight use.

\subsection{When the capacity constraint changes the distribution}
The four-slot star shows an unrelated contrast of 0.546 blocks ([0.429, 0.667]), a full-cohort mean-burden increase of 0.333 blocks, and a Gini change of +0.0184. Its confirmation-incidence contrast is $-0.0625$. These differences combine identity redistribution with a change in the confirmation-count sequence; some generated star chunks exceed the complete four-slot budget. This is why its aggregate distribution can change while the chain's does not. A high initial dependency premium is insufficient to guarantee inclusion when the selection unit does not fit.

Two-slot cases provide more severe capacity stress. The observed Core minus scalar unrelated means remain nonnegative in these cells, but dependent transactions can also fare worse. In the two-slot star, the parent waits 2.900 more blocks on average under Core despite a positive snapshot dependency premium. A selection unit that cannot fit delivers no priority advantage. This is not evidence that the default Mainnet allocator ordinarily excludes whole clusters indefinitely; the maximum default cluster is far smaller than a default full block.

At twelve slots, all twelve transactions are included in the first block under every selector and graph. Fees and utilization agree, confirmation incidence is one, and restricted-burden Gini is zero. Dependencies can still change the snapshot priority premium but generate no delay contrast when all demand fits. This no-congestion null is part of the result.

\subsection{Topology and algorithm comparisons have different meanings}
The released analysis also holds the Core selector fixed and compares each dependent graph with its newly signed independent-graph counterpart. These paired topology contrasts preserve fees and serialization size but do not hold raw identifiers or UTXO funding values fixed. They are distinct from the Core minus scalar contrasts in the main table. For the four-slot star, the Core topology contrast in unrelated burden is $-0.133$ blocks (interval [-0.263, -0.008]), whereas changing from scalar to Core on the star increases it by 0.546 blocks. The counterfactual therefore changes even the sign of the answer. Against the ancestor comparator on the four-slot chain, the unrelated-burden contrast is $-0.042$ blocks (interval [-0.079, -0.008]), with 1,456.775 fewer first-block satoshis on average. A revenue improvement in one selected block is neither a higher total fee payment by the fixed cohort nor evidence of greater social efficiency.
